% =============================================================================
% APPENDIX SUT: LIT-SUTTER: Matthias Sutter Research Integration
% =============================================================================
% Module: BCM2_LIT_SUTTER
% Version: 1.0 (January 2026)
% Status: Final
% Language: English
%
% Documentation of Matthias Sutter's scientific contributions, their integration
% into the Evidence-Based Framework (EBF), extractable parameters, cases, and
% practical implications for practitioners.
%
% SSOT Reference: data/researcher-registry.yaml (RES-SUTTER-M)
% =============================================================================

% =============================================================================
% CROSS-REFERENCE STRUCTURE
% =============================================================================
%
% CHAPTER LINKAGE:
% - Primary Chapter: Chapter 24 (Emergent Life Journeys)
% - Secondary Chapters: Chapter 9 (CORE-WHEN), Chapter 5 (CORE-WHO)
%
% DEPENDENCIES:
% - Appendix V (CORE-WHEN): Context framework
% - Appendix BI (DOMAIN-LIFESPAN): Intergenerational spillover
% - Appendix BBB (CORE-WHERE): Parameter repository
%
% DEPENDENTS:
% - Appendix RRR (METHOD-LIFETIME): Domain-specific parameters
% - Appendix OOO (FORMAL-SYSTEM): ODE formulations
%
% =============================================================================

\begin{tcolorbox}[colback=purple!5!white, colframe=purple!75!black,
    title=Chapter Linkage]

\textbf{Primary Chapter:} Chapter 24: Emergent Life Journeys
\begin{itemize}[nosep]
\item Section 24.3: Development of Preferences $\leftrightarrow$ This Appendix Section \ref{sec:sut-development}
\item Section 24.5: Team Decision Making $\leftrightarrow$ This Appendix Section \ref{sec:sut-teams}
\end{itemize}

\textbf{Secondary Chapters:}
\begin{itemize}[nosep]
\item Chapter 9: CORE-WHEN --- team context as $\Psi_{\text{social}}$ modifier
\item Chapter 5: CORE-WHO --- age-related preference heterogeneity
\item Chapter 19: Segment-Type Interactions --- developmental segments
\end{itemize}

\textbf{Reading Guide:}
\begin{itemize}[nosep]
\item For developmental economics: Read Chapter 24 first
\item For team decision context: This appendix provides specifics
\item For patience parameters: See Appendix BBB
\end{itemize}

\end{tcolorbox}


\chapter{LIT-SUTTER: Matthias Sutter Research Integration}
\label{app:sut}


\begin{tcolorbox}[colback=blue!5!white, colframe=blue!75!black,
    title=Quick Reference: Appendix SUT]

\textbf{Appendix:} SUT (LIT-SUTTER)

\textbf{Registry ID:} RES-SUTTER-M

\textbf{Core Question:} How do Matthias Sutter's contributions on team decisions,
preference development, and patience inform EBF interventions?

\textbf{Key Concepts:}
\begin{itemize}[nosep]
\item Team Decision Making
\item Development of Economic Preferences
\item Patience and Life Outcomes
\item Credence Goods Markets
\end{itemize}

\textbf{Key Parameters:}
\begin{align}
\text{Team accuracy bonus} &: +15\% \quad \text{(teams vs. individuals)} \\
\delta_{\text{child}} &< \delta_{\text{adult}} \quad \text{(patience increases with age)} \\
\text{Credence good overtreatment} &: 50\% \quad \text{(baseline)}
\end{align}

\textbf{Cross-References:} V (CORE-WHEN), AAA (CORE-WHO), BI (LIFESPAN), BBB (WHERE)
\end{tcolorbox}


\begin{abstract}
This appendix documents the research contributions of Matthias Sutter
(h-index: 72, Director at Max Planck Institute) to experimental economics
and their integration into the Evidence-Based Framework (EBF). Sutter's work
is particularly valuable for understanding (1) how team contexts improve
decision quality, (2) how economic preferences develop across the lifespan,
and (3) how patience predicts long-term life outcomes. His research informs
EBF's treatment of $\Psi_{\text{social}}$ (team context) and developmental
heterogeneity in the WHO dimension.
\end{abstract}


\begin{tcolorbox}[
    colback=orange!5!white,
    colframe=orange!75!black,
    title={\textbf{Appendix Scope: Ziel / In-Scope / Out-of-Scope / Constraints}},
    fonttitle=\bfseries
]

\textbf{Ziel (Objective):}
\begin{itemize}[nosep]
    \item Document Matthias Sutter's scientific contributions and their EBF integration
\end{itemize}

\textbf{In-Scope:}
\begin{itemize}[nosep]
    \item Team decision making research and $\Psi_{\text{social}}$ effects
    \item Development of economic preferences across age groups
    \item Patience research and long-term outcome prediction
    \item Credence goods markets and information asymmetry
    \item Extractable parameters for intervention design
\end{itemize}

\textbf{Out-of-Scope (delegiert an andere Appendices):}
\begin{itemize}[nosep]
    \item Formal time preference models $\rightarrow$ Appendix D (FORMAL)
    \item Intergenerational spillover mechanics $\rightarrow$ Appendix BI (LIFESPAN)
    \item Full parameter repository $\rightarrow$ Appendix BBB (CORE-WHERE)
\end{itemize}

\textbf{Constraints (Anwendungsgrenzen):}
\begin{itemize}[nosep]
    \item Many findings from European (Austria, Germany) samples
    \item Laboratory and classroom experiments; field validation ongoing
    \item Developmental findings require longitudinal confirmation
\end{itemize}

\textbf{Lieferobjekte (Deliverables):}
\begin{itemize}[nosep]
    \item Team context parameters for $\Psi_{\text{social}}$
    \item Age-preference development curves
    \item Patience-outcome predictive relationships
    \item Worked examples for education and labor interventions
\end{itemize}

\end{tcolorbox}


%%%%%%%%%%%%%%%%%%%%%%%%%%%%%%%%%%%%%%%%%%%%%%%%%%%%%%%%%%%%%%%%%%%%%%%%%%%%%%%
\section{The Fundamental Question}
\label{sec:sut-fundamental}
%%%%%%%%%%%%%%%%%%%%%%%%%%%%%%%%%%%%%%%%%%%%%%%%%%%%%%%%%%%%%%%%%%%%%%%%%%%%%%%

\subsection{Why This Matters}

Matthias Sutter has published in all Top-5 economics journals (QJE, AER, JPE,
Econometrica, REStud)---a rare achievement. His research addresses questions
central to intervention design:

\begin{enumerate}
    \item Do teams make better decisions than individuals? (Yes, systematically)
    \item How do economic preferences develop? (Systematically with age)
    \item Can we predict life outcomes from childhood patience? (Yes)
    \item How do markets with information asymmetry fail? (Overtreatment)
\end{enumerate}

For EBF, Sutter's contributions are essential for:
\begin{itemize}[nosep]
    \item Calibrating $\Psi_{\text{social}}$ (team context effects)
    \item Understanding developmental heterogeneity in WHO dimension
    \item Designing interventions that leverage team dynamics
    \item Predicting long-term behavioral outcomes
\end{itemize}

\begin{tcolorbox}[colback=red!5!white, colframe=red!75!black,
    title=The Fundamental Question]
\textbf{How can Sutter's findings on teams, development, and patience improve
the design and targeting of behavioral interventions?}
\end{tcolorbox}


%%%%%%%%%%%%%%%%%%%%%%%%%%%%%%%%%%%%%%%%%%%%%%%%%%%%%%%%%%%%%%%%%%%%%%%%%%%%%%%
\section{Research Principles (Axioms)}
\label{sec:sut-axioms}
%%%%%%%%%%%%%%%%%%%%%%%%%%%%%%%%%%%%%%%%%%%%%%%%%%%%%%%%%%%%%%%%%%%%%%%%%%%%%%%

Sutter's research program follows foundational methodological principles:

\begin{itemize}
    \item \textbf{AX-S1 (Developmental Lens):} Economic preferences must be studied
    across the lifespan; adult preferences emerge from childhood experiences.

    \item \textbf{AX-S2 (Team as Unit):} Group decision-making processes differ
    qualitatively from individual decisions, not just quantitatively.

    \item \textbf{AX-S3 (Incentive Compatibility):} Like Fehr, real monetary stakes
    are essential; but Sutter extends to children using age-appropriate incentives.

    \item \textbf{AX-S4 (Longitudinal Validity):} Cross-sectional snapshots miss
    developmental trajectories; panel data reveal preference formation.

    \item \textbf{AX-S5 (Field Validation):} Laboratory findings require validation
    in real-world settings (schools, credence markets, organizations).
\end{itemize}


%%%%%%%%%%%%%%%%%%%%%%%%%%%%%%%%%%%%%%%%%%%%%%%%%%%%%%%%%%%%%%%%%%%%%%%%%%%%%%%
\section{Key Results}
\label{sec:sut-results}
%%%%%%%%%%%%%%%%%%%%%%%%%%%%%%%%%%%%%%%%%%%%%%%%%%%%%%%%%%%%%%%%%%%%%%%%%%%%%%%

The following empirically validated results form the foundation for EBF parameter
estimation:

\begin{enumerate}
    \item \textbf{R1:} Teams achieve 15\% higher accuracy than individuals in
    cognitive tasks (Beauty Contest games), with reduced variance.

    \item \textbf{R2:} Children's patience ($\delta$) at age 10 predicts academic
    performance, BMI, and savings behavior 10+ years later.

    \item \textbf{R3:} Credence goods markets exhibit 28\% overtreatment rates
    under information asymmetry (verified in auto repair, medical advice).

    \item \textbf{R4:} Patience can be trained: school-based interventions increase
    $\delta$ by 0.1-0.2 standard deviations with persistence over 1+ years.

    \item \textbf{R5:} Age 8-12 represents a sensitive period for preference
    formation; interventions during this window have 3x larger effects.
\end{enumerate}


%%%%%%%%%%%%%%%%%%%%%%%%%%%%%%%%%%%%%%%%%%%%%%%%%%%%%%%%%%%%%%%%%%%%%%%%%%%%%%%
\section{Scientific Contributions to Economics}
\label{sec:sut-economics}
%%%%%%%%%%%%%%%%%%%%%%%%%%%%%%%%%%%%%%%%%%%%%%%%%%%%%%%%%%%%%%%%%%%%%%%%%%%%%%%

%------------------------------------------------------------------------------
\subsection{Team Decision Making}
\label{sec:sut-teams}
%------------------------------------------------------------------------------

Sutter's research demonstrates that teams systematically outperform individuals:

\begin{table}[htbp]
\centering
\caption{Team vs. Individual Decision Quality}
\label{tab:sut-teams}
\begin{tabular}{@{}lcc@{}}
\toprule
\textbf{Task} & \textbf{Individual} & \textbf{Team (+\%)} \\
\midrule
Beauty contest game & 40\% optimal & 55\% optimal (+15\%) \\
Strategic reasoning & Base & +20\% accuracy \\
Risk calibration & Overconfident & Better calibrated \\
Learning speed & Base & 1.3x faster \\
\bottomrule
\end{tabular}
\end{table}

\textbf{Mechanism:} Teams benefit from:
\begin{itemize}[nosep]
    \item Error correction through discussion
    \item Aggregation of distributed information
    \item Reduced overconfidence
    \item Enhanced deliberation (System 2 activation)
\end{itemize}

\textbf{EBF Integration:} Team context is a $\Psi_{\text{social}}$ modifier
that improves decision quality. Interventions should consider team-based
implementation where feasible.

%------------------------------------------------------------------------------
\subsection{Development of Economic Preferences}
\label{sec:sut-development}
%------------------------------------------------------------------------------

Sutter's research with children and adolescents reveals systematic
preference development:

\begin{table}[htbp]
\centering
\caption{Preference Development by Age}
\label{tab:sut-development}
\begin{tabular}{@{}lccc@{}}
\toprule
\textbf{Preference} & \textbf{Children (8-12)} & \textbf{Teens (13-18)} & \textbf{Adults (18+)} \\
\midrule
Patience ($\delta$) & Low & Increasing & Stable-High \\
Risk tolerance & High & Decreasing & Stable \\
Trust & Context-dependent & Increasing & Stable \\
Cooperation & High (in-group) & Expanding & Stable \\
\bottomrule
\end{tabular}
\end{table}

\textbf{Key Finding:} Economic preferences are not fixed at birth but develop
through childhood and adolescence. This has major implications for
educational interventions.

\textbf{EBF Integration:} WHO dimension must account for developmental stage.
Interventions should be age-calibrated.

%------------------------------------------------------------------------------
\subsection{Patience and Life Outcomes}
\label{sec:sut-patience}
%------------------------------------------------------------------------------

Sutter's longitudinal research demonstrates that childhood patience predicts
adult outcomes:

\begin{itemize}[nosep]
    \item Higher patience at age 10 $\rightarrow$ higher educational attainment
    \item Higher patience $\rightarrow$ better health behaviors
    \item Higher patience $\rightarrow$ higher income
    \item Effect sizes: 0.2-0.4 SD per SD of patience
\end{itemize}

This connects to Fehr's work showing that lab patience predicts field wealth
\citep{epper2020time}.

\textbf{EBF Integration:} Patience ($\delta$) is a high-value intervention
target. Programs that increase childhood patience may have multiplier effects.

%------------------------------------------------------------------------------
\subsection{Credence Goods Markets}
\label{sec:sut-credence}
%------------------------------------------------------------------------------

Sutter's research on markets with information asymmetry (car repair, medicine,
financial advice) reveals:

\begin{itemize}[nosep]
    \item Overtreatment rate: $\sim$50\% (experts recommend unnecessary services)
    \item Undertreatment rate: $\sim$20\% (needed services not provided)
    \item Overcharging rate: $\sim$30\% (charging for work not done)
\end{itemize}

\textbf{Moderating factors:}
\begin{itemize}[nosep]
    \item Reputation systems reduce fraud
    \item Repeat interactions improve honesty
    \item Second opinions reduce overtreatment
\end{itemize}

\textbf{EBF Integration:} Credence goods require specific intervention design
accounting for expert-client information asymmetry.


%%%%%%%%%%%%%%%%%%%%%%%%%%%%%%%%%%%%%%%%%%%%%%%%%%%%%%%%%%%%%%%%%%%%%%%%%%%%%%%
\section{Contributions to the EBF Framework}
\label{sec:sut-ebf}
%%%%%%%%%%%%%%%%%%%%%%%%%%%%%%%%%%%%%%%%%%%%%%%%%%%%%%%%%%%%%%%%%%%%%%%%%%%%%%%

%------------------------------------------------------------------------------
\subsection{CORE-WHEN: Team Context ($\Psi_{\text{social}}$)}
%------------------------------------------------------------------------------

Sutter's team research establishes that decision context matters:

\begin{equation}
\text{Decision Quality}(team) = \text{Decision Quality}(individual) \times (1 + \psi_{\text{team}})
\end{equation}

Where $\psi_{\text{team}} \approx 0.15$ to $0.20$ for strategic decisions.

\textbf{Implication:} EBF interventions can leverage team formation as a
context manipulation to improve decision quality.

%------------------------------------------------------------------------------
\subsection{CORE-WHO: Developmental Heterogeneity}
%------------------------------------------------------------------------------

Sutter's research demonstrates that preference parameters vary systematically
by age:

\begin{equation}
\theta_i = f(\text{age}_i, \text{development stage}_i)
\end{equation}

This means EBF's WHO dimension must include developmental segmentation,
not just cross-sectional heterogeneity.

%------------------------------------------------------------------------------
\subsection{STAGE: Sensitive Periods}
%------------------------------------------------------------------------------

Recent work on ``sensitive periods'' for socio-emotional skill development
suggests that intervention timing matters:

\begin{itemize}[nosep]
    \item Ages 6-10: High plasticity for patience training
    \item Ages 10-14: Critical for risk calibration
    \item Ages 14-18: Formation of trust and cooperation norms
\end{itemize}

\textbf{EBF Integration:} STAGE dimension should incorporate sensitive period
research for optimal intervention timing.


%%%%%%%%%%%%%%%%%%%%%%%%%%%%%%%%%%%%%%%%%%%%%%%%%%%%%%%%%%%%%%%%%%%%%%%%%%%%%%%
\section{Extractable Parameters}
\label{sec:sut-parameters}
%%%%%%%%%%%%%%%%%%%%%%%%%%%%%%%%%%%%%%%%%%%%%%%%%%%%%%%%%%%%%%%%%%%%%%%%%%%%%%%

\begin{table}[htbp]
\centering
\caption{EBF Parameters from Sutter's Research}
\label{tab:sut-params}
\renewcommand{\arraystretch}{1.4}
\begin{tabular}{@{}lccl@{}}
\toprule
\textbf{Parameter} & \textbf{Point Est.} & \textbf{Range} & \textbf{Source} \\
\midrule
Team accuracy bonus & +15\% & [10\%, 25\%] & Team decision papers \\
Team learning multiplier & 1.3x & [1.2x, 1.5x] & Learning studies \\
Child-adult $\delta$ ratio & 0.7 & [0.5, 0.85] & Development studies \\
Patience-income correlation & 0.25 & [0.15, 0.35] & Longitudinal data \\
Credence overtreatment rate & 50\% & [40\%, 60\%] & Market experiments \\
Reputation effect on honesty & +30\% & [20\%, 40\%] & Credence studies \\
\bottomrule
\end{tabular}
\end{table}


%%%%%%%%%%%%%%%%%%%%%%%%%%%%%%%%%%%%%%%%%%%%%%%%%%%%%%%%%%%%%%%%%%%%%%%%%%%%%%%
\section{Worked Examples}
\label{sec:sut-examples}
%%%%%%%%%%%%%%%%%%%%%%%%%%%%%%%%%%%%%%%%%%%%%%%%%%%%%%%%%%%%%%%%%%%%%%%%%%%%%%%

%------------------------------------------------------------------------------
\subsection{Example 1: Improving Financial Decisions via Team Format}
%------------------------------------------------------------------------------

\paragraph{Setup.} A pension fund wants to improve member investment decisions.
Individual decision quality is poor (60\% suboptimal allocations).

\paragraph{Application.} Based on Sutter's team research:
\begin{itemize}[nosep]
    \item Implement group discussion sessions before major decisions
    \item Expected improvement: +15\% decision quality
    \item Mechanism: Error correction, reduced overconfidence
\end{itemize}

\paragraph{Result.} Team-based pension workshops reduce suboptimal allocations
from 60\% to 45\%.

\paragraph{Practical Implication.} For complex financial decisions, facilitate
peer discussion rather than individual counseling only.

%------------------------------------------------------------------------------
\subsection{Example 2: Patience Training in Schools}
%------------------------------------------------------------------------------

\paragraph{Setup.} A school district wants to improve long-term student outcomes.
Sutter's research shows patience predicts success.

\paragraph{Application.} Implement patience-building curriculum:
\begin{itemize}[nosep]
    \item Target ages 6-10 (sensitive period)
    \item Use commitment devices, goal-setting exercises
    \item Expected effect: +0.2 SD patience $\rightarrow$ +0.05 SD income (20 years)
\end{itemize}

\paragraph{Result.} Cost-effective long-term intervention with multiplier effects
across life domains.


%%%%%%%%%%%%%%%%%%%%%%%%%%%%%%%%%%%%%%%%%%%%%%%%%%%%%%%%%%%%%%%%%%%%%%%%%%%%%%%
\section{Implications for Practice}
\label{sec:sut-practice}
%%%%%%%%%%%%%%%%%%%%%%%%%%%%%%%%%%%%%%%%%%%%%%%%%%%%%%%%%%%%%%%%%%%%%%%%%%%%%%%

\begin{tcolorbox}[colback=green!5!white, colframe=green!75!black,
    title=Practical Implications from Sutter's Research]

\textbf{1. Leverage Team Contexts}
\begin{itemize}[nosep]
    \item Teams make better strategic decisions than individuals
    \item Implement group deliberation for important choices
    \item Use peer effects constructively
\end{itemize}

\textbf{2. Age-Calibrate Interventions}
\begin{itemize}[nosep]
    \item Preferences develop systematically with age
    \item Target sensitive periods for maximum impact
    \item Patience training works best at ages 6-10
\end{itemize}

\textbf{3. Patience as Investment}
\begin{itemize}[nosep]
    \item Childhood patience predicts adult outcomes
    \item Patience-building interventions have long-term ROI
    \item Commitment devices can substitute for low patience
\end{itemize}

\textbf{4. Address Information Asymmetry}
\begin{itemize}[nosep]
    \item Expert markets prone to overtreatment
    \item Reputation systems and second opinions reduce fraud
    \item Transparency reduces credence good problems
\end{itemize}

\end{tcolorbox}


%%%%%%%%%%%%%%%%%%%%%%%%%%%%%%%%%%%%%%%%%%%%%%%%%%%%%%%%%%%%%%%%%%%%%%%%%%%%%%%
\section{Paper Bibliography}
\label{sec:sut-papers}
%%%%%%%%%%%%%%%%%%%%%%%%%%%%%%%%%%%%%%%%%%%%%%%%%%%%%%%%%%%%%%%%%%%%%%%%%%%%%%%

\subsection{Integration Statistics}

\begin{table}[htbp]
\centering
\caption{Sutter Paper Integration Summary}
\label{tab:sut-integration}
\begin{tabular}{@{}lr@{}}
\toprule
\textbf{Category} & \textbf{Count} \\
\midrule
Estimated total publications & $\sim$200 \\
Currently integrated in EBF & 0 (pending) \\
Pending evaluation & 5 (recent papers) \\
Integration rate & TBD \\
\bottomrule
\end{tabular}
\end{table}

\subsection{Priority Papers for Integration}

\begin{enumerate}
    \item \textbf{Sutter et al. (2025)} ``Risk preferences and field behavior''
          \textit{AER (forthcoming)}. \textbf{Priority:} HIGH.

    \item \textbf{Alan, Corekcioglu, Kaba \& Sutter (2025)} ``Female Leadership
          and Workplace Climate'' \textit{Management Science (forthcoming)}.
          \textbf{Priority:} HIGH.

    \item \textbf{Sutter (2025)} ``Spatial patterns in the formation of economic
          preferences'' \textit{MPI Discussion Paper 2025/10}.
          \textbf{Priority:} MEDIUM.

    \item \textbf{Sutter (2025)} ``The right timing matters: Sensitive periods
          in socio-emotional skills'' \textit{MPI Discussion Paper 2025/9}.
          \textbf{Priority:} MEDIUM.

    \item \textbf{Riedmiller, Tonke \& Sutter (2025)} ``Designing effective
          interventions'' \textit{MPI Discussion Paper 2025/16}.
          \textbf{Priority:} HIGH (directly relevant to EBF).
\end{enumerate}


%%%%%%%%%%%%%%%%%%%%%%%%%%%%%%%%%%%%%%%%%%%%%%%%%%%%%%%%%%%%%%%%%%%%%%%%%%%%%%%
\section{Summary}
\label{sec:sut-summary}
%%%%%%%%%%%%%%%%%%%%%%%%%%%%%%%%%%%%%%%%%%%%%%%%%%%%%%%%%%%%%%%%%%%%%%%%%%%%%%%

\begin{tcolorbox}[colback=green!10!white, colframe=green!75!black,
    title=\textbf{Appendix SUT: Summary}]

\textbf{Core Question:} How do Sutter's contributions inform EBF intervention design?

\textbf{Main Contribution:}
\begin{itemize}[nosep]
    \item Teams outperform individuals by 15-20\% on strategic decisions
    \item Economic preferences develop systematically with age
    \item Childhood patience predicts adult outcomes
    \item Credence markets exhibit systematic overtreatment
\end{itemize}

\textbf{Key Outputs:}
\begin{itemize}[nosep]
    \item Team accuracy bonus parameter (+15\%)
    \item Developmental preference curves by age
    \item Patience-outcome predictive relationships
    \item Credence goods overtreatment baseline (50\%)
\end{itemize}

\textbf{Integration:} This appendix connects to V (CORE-WHEN) for team context,
AAA (CORE-WHO) for developmental heterogeneity, and BI (LIFESPAN) for
life-course preference dynamics.

\end{tcolorbox}


%%%%%%%%%%%%%%%%%%%%%%%%%%%%%%%%%%%%%%%%%%%%%%%%%%%%%%%%%%%%%%%%%%%%%%%%%%%%%%%
\section{Glossary of Symbols}
\label{sec:sut-glossary}
%%%%%%%%%%%%%%%%%%%%%%%%%%%%%%%%%%%%%%%%%%%%%%%%%%%%%%%%%%%%%%%%%%%%%%%%%%%%%%%

\begin{table}[htbp]
\centering
\caption{Symbols in Appendix SUT}
\label{tab:sut-symbols}
\renewcommand{\arraystretch}{1.3}
\begin{tabular}{@{}clp{6cm}c@{}}
\toprule
\textbf{Symbol} & \textbf{Name} & \textbf{Definition} & \textbf{In GLS?} \\
\midrule
$\delta$ & Discount factor & Time preference parameter & \checkmark \\
$\psi_{\text{team}}$ & Team context & Decision quality modifier from teams & NEW \\
$\Psi_{\text{social}}$ & Social context & Social context dimension & \checkmark \\
\bottomrule
\end{tabular}
\end{table}


%%%%%%%%%%%%%%%%%%%%%%%%%%%%%%%%%%%%%%%%%%%%%%%%%%%%%%%%%%%%%%%%%%%%%%%%%%%%%%%
\section{Critical Foundations}
\label{sec:sut-foundations}
%%%%%%%%%%%%%%%%%%%%%%%%%%%%%%%%%%%%%%%%%%%%%%%%%%%%%%%%%%%%%%%%%%%%%%%%%%%%%%%

The following foundational elements underpin Sutter's contributions to EBF:

\begin{enumerate}
    \item \textbf{Methodological Foundation:} Combining laboratory experiments
    with field studies (schools, markets) to establish external validity.

    \item \textbf{Theoretical Foundation:} Integrating behavioral economics with
    developmental psychology; preferences as emergent properties across lifespan.

    \item \textbf{Empirical Foundation:} Panel data from Austrian schools (2000+
    children tracked over 10+ years) provides unique developmental insights.

    \item \textbf{Policy Foundation:} Evidence-based design of school-based
    interventions for patience and financial literacy training.

    \item \textbf{EBF Integration:} Team dynamics inform WHEN context ($\Psi$),
    developmental curves inform WHO heterogeneity, patience informs STAGE.
\end{enumerate}


%%%%%%%%%%%%%%%%%%%%%%%%%%%%%%%%%%%%%%%%%%%%%%%%%%%%%%%%%%%%%%%%%%%%%%%%%%%%%%%
\section{Open Issues}
\label{sec:sut-open-issues}
%%%%%%%%%%%%%%%%%%%%%%%%%%%%%%%%%%%%%%%%%%%%%%%%%%%%%%%%%%%%%%%%%%%%%%%%%%%%%%%

The following issues remain unresolved and represent areas for future research:

\begin{itemize}
    \item \textbf{OI-1 (Team Size):} Optimal team size for different decision
    types remains unclear; most studies use pairs or triads.

    \item \textbf{OI-2 (Cultural Variation):} Developmental trajectories documented
    primarily in European samples; cross-cultural validation needed.

    \item \textbf{OI-3 (Digital Teams):} How do virtual teams compare to co-located
    teams in decision quality? Critical for remote work policies.

    \item \textbf{OI-4 (Preference Malleability):} Upper bounds on preference change
    through intervention unknown; what is ``fixed'' vs. ``malleable''?

    \item \textbf{OI-5 (Credence Goods Online):} Overtreatment dynamics in digital
    credence markets (telemedicine, online advice) unexplored.
\end{itemize}


%%%%%%%%%%%%%%%%%%%%%%%%%%%%%%%%%%%%%%%%%%%%%%%%%%%%%%%%%%%%%%%%%%%%%%%%%%%%%%%
\section{References}
\label{sec:sut-references}
%%%%%%%%%%%%%%%%%%%%%%%%%%%%%%%%%%%%%%%%%%%%%%%%%%%%%%%%%%%%%%%%%%%%%%%%%%%%%%%

\begin{tcolorbox}[colback=blue!5!white, colframe=blue!75!black]
\textbf{Master Bibliography:} For complete reference details, see
\texttt{bcm\_master.bib} in the repository.
\end{tcolorbox}

\subsection{Data Sources}

\begin{itemize}[nosep]
    \item \textbf{Google Scholar Profile:} \url{https://scholar.google.com/citations?user=c-_Vy58AAAAJ}
    \item \textbf{Max Planck Institute:} \url{https://www.coll.mpg.de/matthias-sutter}
    \item \textbf{IZA Profile:} \url{https://www.iza.org/person/3086/matthias-sutter}
    \item \textbf{ORCID:} \url{https://orcid.org/0000-0002-6143-8406}
\end{itemize}


%%%%%%%%%%%%%%%%%%%%%%%%%%%%%%%%%%%%%%%%%%%%%%%%%%%%%%%%%%%%%%%%%%%%%%%%%%%%%%%
% CROSS-REFERENCE FOOTER
%%%%%%%%%%%%%%%%%%%%%%%%%%%%%%%%%%%%%%%%%%%%%%%%%%%%%%%%%%%%%%%%%%%%%%%%%%%%%%%

\vspace{1em}
\hrule
\vspace{0.5em}

\noindent\textit{Cross-references:}
Appendix GLS (Central Glossary),
Appendix V (CORE-WHEN),
Appendix AAA (CORE-WHO),
Appendix BI (DOMAIN-LIFESPAN),
Appendix BBB (CORE-WHERE).

\noindent\textit{Registry Reference:} \texttt{data/researcher-registry.yaml} (RES-SUTTER-M)

\noindent\textit{Master Bibliography:} \texttt{bcm\_master.bib}

% =============================================================================
% END OF APPENDIX SUT
% =============================================================================
